\section{张量并行：按宽度切分}

数据并行解决的是“模型放得下吗”，张量并行解决的是更难的问题：\textbf{如果单层矩阵都太大，怎么把一层内部再切开。} 它的思路是切 hidden dimension，让每个 rank 只持有一部分参数和一部分中间激活。

\begin{figure}[H]
\centering
\begin{tikzpicture}[every node/.style={font=\small, align=center}]
\node[draw, rounded corners, fill=orange!10, minimum width=3cm, minimum height=0.9cm] (full) at (0,0) {完整层\\$W \in \mathbb{R}^{d \times d}$};
\node[draw, rounded corners, fill=orange!18, minimum width=3cm, minimum height=0.9cm] (part1) at (4,0) {rank 0\\$W_0$};
\node[draw, rounded corners, fill=orange!18, minimum width=3cm, minimum height=0.9cm] (part2) at (8,0) {rank 1\\$W_1$};
\node[draw, rounded corners, fill=orange!18, minimum width=3cm, minimum height=0.9cm] (part3) at (12,0) {rank 2\\$W_2$};
\draw[->, thick] (full) -- (part1);
\draw[->, thick] (full) -- (part2);
\draw[->, thick] (full) -- (part3);
\end{tikzpicture}
\caption{张量并行的核心：把一层里的宽度切成多份}
\end{figure}

\subsection{工作方式}

在 lecture 的简化 MLP 里，每层都是矩阵乘法加非线性。张量并行的做法是：每个 rank 只算自己的参数切片，算完之后再把激活拼回完整形状。于是，\textbf{参数省了，通信却变多了}。

\begin{figure}[H]
\centering
\begin{tikzpicture}[every node/.style={font=\small, align=center}]
\node[draw, rounded corners, fill=blue!10, minimum width=3.1cm, minimum height=0.9cm] (local) at (0,0) {本地切片\\$x @ W_i$};
\node[draw, rounded corners, fill=blue!18, minimum width=3.1cm, minimum height=0.9cm] (ag) at (4.2,0) {all-gather\\收集所有切片};
\node[draw, rounded corners, fill=blue!26, minimum width=3.1cm, minimum height=0.9cm] (cat) at (8.4,0) {concat\\恢复完整激活};
\node[draw, rounded corners, fill=blue!34, minimum width=3.1cm, minimum height=0.9cm] (next) at (12.6,0) {下一层\\继续切分};
\draw[->, thick] (local) -- (ag);
\draw[->, thick] (ag) -- (cat);
\draw[->, thick] (cat) -- (next);
\end{tikzpicture}
\caption{张量并行的前向过程：每层都需要把部分结果重新拼回去}
\end{figure}

\begin{importantbox}{为什么张量并行要求更快的互联}
因为它往往是“每层都要通信”。如果互联慢，模型虽然被切开了，但中间激活会反复来回搬，最后把计算收益全吃掉。
\end{importantbox}

\begin{figure}[H]
\centering
\begin{tikzpicture}[every node/.style={font=\small, align=center}]
\node[draw, rounded corners, fill=yellow!10, minimum width=3.6cm, minimum height=0.9cm] (mem) at (0,0) {每卡参数更少\\更容易放下大模型};
\node[draw, rounded corners, fill=yellow!16, minimum width=3.6cm, minimum height=0.9cm] (comm) at (4.8,0) {每层都要 all-gather\\通信频率很高};
\node[draw, rounded corners, fill=yellow!22, minimum width=3.6cm, minimum height=0.9cm] (hw) at (9.6,0) {需要高带宽互联\\否则通信拖死训练};
\end{tikzpicture}
\caption{张量并行的三角权衡：更省显存，但更吃互联}
\end{figure}

\begin{warningbox}{张量并行的脆弱点}
如果 interconnect 不够快，张量并行会比数据并行更容易被通信拖垮。它适合”单层就已经很大”的模型，但不适合拿来粗暴替代所有其他并行策略。
\end{warningbox}

\begin{warningbox}{张量并行要求低延迟互联——只适合 NVLink 范围}
张量并行的通信发生在\textbf{每一层}的前向和反向中，而不是像 DDP 那样每步只做一次。对于一个 40 层的 transformer，张量并行意味着每步至少 $40 \times 2 = 80$ 次 all-gather 或 reduce-scatter（前向 + 反向各一次）。如果互联延迟从 NVLink 的 $\sim$1 $\mu$s 变为 InfiniBand 的 $\sim$5 $\mu$s，80 次调用的累积延迟从 80 $\mu$s 变为 400 $\mu$s——看起来不大，但加上带宽瓶颈，影响会成倍放大。

\textbf{实践准则}：张量并行只在 NVLink 可达的 GPU 之间使用（通常是同一节点的 8 张卡）。跨节点的并行需求应由数据并行或流水线并行承担。
\end{warningbox}

\subsection{张量并行 vs 数据并行：通信量对比}

把 TP 和 DP 的通信开销放在同一个坐标系里看，能更直观地理解两者的权衡。假设模型有 $L$ 层，每层参数大小为 $W$，激活大小为 $A$：

\begin{table}[H]
\centering
\small
\resizebox{\textwidth}{!}{%
\begin{tabular}{p{2.8cm}p{3.5cm}p{3.5cm}p{3.5cm}}
\toprule
\textbf{维度} & \textbf{数据并行 (DDP)} & \textbf{张量并行 (TP)} & \textbf{比较} \\
\midrule
通信频率 & 每步 1 次 & 每层 2 次（前向+反向） & TP 频率高 $\sim 2L$ 倍 \\
每次通信量 & $2 \cdot \frac{p-1}{p} \cdot LW$ & $2 \cdot \frac{p-1}{p} \cdot A$ & 取决于 $LW$ vs $A$ \\
总通信量/步 & $2 \cdot \frac{p-1}{p} \cdot LW$ & $2L \cdot 2 \cdot \frac{p-1}{p} \cdot A$ & DP 通信与参数成正比 \\
通信类型 & all-reduce（梯度） & all-gather + reduce-scatter & TP 更敏感于延迟 \\
互联要求 & InfiniBand 可接受 & 必须 NVLink & TP 对硬件更挑剔 \\
\bottomrule
\end{tabular}
}
\caption{数据并行 vs 张量并行通信量对比}
\end{table}

\begin{knowledgebox}{一个具体的数字对比}
以 Llama 2 13B（40 层，$d=5120$，$B=128$，$p=4$，fp16）为例：
\begin{itemize}
    \item \textbf{DDP 每步总通信量}：$2 \times \frac{3}{4} \times 26\,\text{GB} = 39\,\text{GB}$，但只需 1 次 all-reduce。
    \item \textbf{TP 每步总通信量}：每层前向 all-gather 激活 $\approx 128 \times 5120 \times 2 = 1.25\,\text{MB}$，40 层前向+反向共 $\approx 80 \times 1.25\,\text{MB} \times \frac{3}{4} \times 2 \approx 150\,\text{MB}$。
    \item TP 的总字节数远小于 DDP，但它需要 80 次独立的集合操作，每次都有启动延迟。在 NVLink 上这不是问题，在跨节点互联上就会成为致命瓶颈。
\end{itemize}
\end{knowledgebox}

\subsection{为什么张量并行的反向传播更难}

前向里，我们还能把一层的输出拆开并在必要时拼回去；但在反向里，梯度流向会反过来，意味着你常常既要知道“每个 shard 自己贡献了什么”，也要知道“所有 shard 一起到底形成了什么”。这就是为什么张量并行经常会配合 \texttt{all\_gather}、\texttt{reduce\_scatter} 或者等价的分片同步原语。

\begin{importantbox}{张量并行的关键不是切片，而是配套的梯度路径}
只把矩阵切开还不够。前向、反向、参数更新必须共享同一套分片约定，否则每个 rank 看到的只是局部真相，拼不回正确的整体梯度。
\end{importantbox}

\begin{figure}[H]
\centering
\begin{tikzpicture}[every node/.style={font=\small, align=center}]
\node[draw, rounded corners, fill=orange!10, minimum width=3.4cm, minimum height=0.9cm] (fw) at (0,0) {前向\\局部矩阵乘};
\node[draw, rounded corners, fill=orange!18, minimum width=3.4cm, minimum height=0.9cm] (sync) at (4.5,0) {同步\\拼接或分片归约};
\node[draw, rounded corners, fill=orange!26, minimum width=3.4cm, minimum height=0.9cm] (bw) at (9.0,0) {反向\\梯度沿 shard 回流};
\draw[->, thick] (fw) -- (sync);
\draw[->, thick] (sync) -- (bw);
\end{tikzpicture}
\caption{张量并行的前后向必须配套设计，否则局部计算无法恢复全局正确性}
\end{figure}

\begin{lstlisting}[caption={张量并行的极简理解},label={lst:tp}]
x = full_batch
for W_local in shard_of_each_layer:
    y_local = gelu(x @ W_local)
    all_parts = all_gather(y_local)
    x = cat(all_parts, dim=-1)
\end{lstlisting}

\begin{knowledgebox}{张量并行的核心直觉}
它不是把模型“切碎”就结束，而是每一层之后都要恢复形状。换句话说，张量并行把压力从“显存”转移到了“通信链路”。
\end{knowledgebox}

\subsection{一个张量并行的数值化例子}

把上面的直觉换成数字会更清楚。假设输入激活大小是 $B \times d$，其中 batch $B=128$，hidden width $d=1024$，采用 fp16，那么单个激活张量的大小大约是
\[
128 \times 1024 \times 2 \text{ bytes} = 256 \text{ KiB}.
\]
如果把 hidden width 切成 4 份，那么每个 rank 先算自己的 $\frac{1}{4}$ 激活，再通过 all-gather 恢复完整张量。按 ring 近似，每个 rank 一层的通信量大约是
\[
2\frac{p-1}{p}A = 1.5A \approx 384 \text{ KiB},
\]
其中 $A$ 是完整激活大小。这个数看上去不大，但注意它是 \textbf{每层都要发生}。

\begin{table}[H]
\centering
\small
\begin{tabular}{p{3.1cm}p{3.8cm}p{4.2cm}}
\toprule
\textbf{量} & \textbf{数值例子} & \textbf{解释} \\
\midrule
batch $B$ & 128 & 讲里常见的小批量起点 \\
hidden $d$ & 1024 & 一个不算夸张的层宽 \\
激活大小 $A$ & 256 KiB & 单层完整输出 \\
4 卡 all-gather 代价 & 384 KiB / layer / rank & 频繁发生，所以要快链路 \\
\bottomrule
\end{tabular}
\caption{张量并行的数字直觉}
\end{table}

\begin{knowledgebox}{为什么张量并行更像“层内系统优化”}
因为它让通信频率提高到“按层计”，于是单层算得再快，也必须保证通信不会把收益抵消。这和 DDP “按步同步” 的节奏完全不同。
\end{knowledgebox}

\begin{figure}[H]
\centering
\begin{tikzpicture}[every node/.style={font=\small, align=center}]
\node[draw, rounded corners, fill=purple!10, minimum width=3.8cm, minimum height=0.9cm] (shard) at (0,0) {shard width\\切 hidden dimension};
\node[draw, rounded corners, fill=purple!18, minimum width=3.8cm, minimum height=0.9cm] (comm) at (5.0,0) {per-layer comm\\每层都交换激活};
\node[draw, rounded corners, fill=purple!26, minimum width=3.8cm, minimum height=0.9cm] (need) at (10.0,0) {fast fabric\\需要高速互联};
\end{tikzpicture}
\caption{张量并行的三步逻辑：切宽度、传激活、靠高速互联兜底}
\end{figure}

